% from ega4-2-fr.tex
\oldpage[IV-2]{155}

tel que $\mathcal F|_{u^{-1}(U)}$ soit plat sur $U$. Si $Y$ est alors le
sous-préschéma réduit de $S$ ayant pour espace sous-jacent $S-U$, il y a par
hypothèse une partition $(Y_i)$ de $Y$ en ensembles localement fermés dans
$Y$ (donc dans $S$) tels que
$\mathcal F_i=\mathcal F\otimes_{\mathcal O_S}\mathcal O_{Y_i}$ soit plat
sur $Y_i$ pour tout $i$; il est clair que les $Y_i$ et $U$ forment une
partition répondant à la question.
\end{proof}

\subsection{Dimension et profondeur d'un Module normalement plat le long d'un
sous-préschéma fermé}
\label{subsection:IV.6.10}

\begin{env}[6.10.1]
\EGAReaderTarget{ega4.c04.statement.6-10-1}{IV.2.s6.10.1-fr}
Soient $X$ un préschéma localement noethérien, $\mathcal J$ un Idéal
quasi-cohérent de $\mathcal O_X$, $Y$ le sous-préschéma fermé de $X$ défini
par $\mathcal J$, $j:Y\to X$ l'injection canonique. Pour tout entier $k\geq0$,
le $\mathcal O_X$-Module
$\mathcal J^k\mathcal F/\mathcal J^{k+1}\mathcal F$ est annulé par
$\mathcal J$, donc de la forme $j_*(\mathcal G_k)$, où
\[
\mathcal G_k=j^*(\mathcal J^k\mathcal F/\mathcal J^{k+1}\mathcal F)
=j^*(\mathcal J^k\mathcal F)
\]
est un $\mathcal O_Y$-Module cohérent. Par abus de langage, nous noterons
$\operatorname{gr}_{\mathcal J}^{\bullet}(\mathcal F)$ le
$\mathcal O_Y$-Module gradué égal à la somme directe
\[
\bigoplus_{k=0}^{\infty}\mathcal G_k
=j^*\left(\bigoplus_{k=0}^{\infty}
\mathcal J^k\mathcal F/\mathcal J^{k+1}\mathcal F\right);
\]
en particulier, on a
$\operatorname{gr}_{\mathcal J}^{0}(\mathcal F)=\mathcal G_0
=\mathcal F\otimes_{\mathcal O_X}\mathcal O_Y=j^*(\mathcal F)$. Nous dirons
(avec Hironaka) que $\mathcal F$ est normalement plat le long de $Y$ si
$\operatorname{gr}_{\mathcal J}^{\bullet}(\mathcal F)$ est un
$\mathcal O_Y$-Module plat. Il revient au même ((2.1.12) et
$(0_{\mathrm I},6.1.2)$) de dire que chacun des $\mathcal O_Y$-Modules
$\mathcal G_k=j^*(\mathcal J^k\mathcal F)$ est localement libre.
\end{env}

\begin{proposition}[6.10.2]
\label{IV.2.s6.10.2-fr}
Soient $X$ un préschéma localement noethérien, $Y$ un sous-préschéma fermé
intègre de $X$. Pour tout $\mathcal O_X$-Module cohérent $\mathcal F$, il
existe un ouvert $U$ de $X$ tel que $Y\cap U\ne\varnothing$ et que
$\mathcal F|U$ soit normalement plat le long de $Y\cap U$.
\end{proposition}

\begin{proof}
Soit en effet $\mathcal J$ l'Idéal cohérent de $\mathcal O_X$ définissant
$Y$; la $\mathcal O_Y$-Algèbre
$\mathcal B=\operatorname{gr}_{\mathcal J}^{\bullet}(\mathcal O_X)$ est
quasi-cohérente et de type fini car elle est engendrée par
$\operatorname{gr}_{\mathcal J}^{1}(\mathcal O_X)$, image réciproque de
$\mathcal J/\mathcal J^2$; comme
$\operatorname{gr}_{\mathcal J}^{\bullet}(\mathcal F)$ est un
$\mathcal B$-module quasi-cohérent engendré par
$\operatorname{gr}_{\mathcal J}^{0}(\mathcal F)$, c'est un
$\mathcal B$-Module de type fini. Si l'on pose
$Z=\operatorname{Spec}(\mathcal B)$, le morphisme structural $u:Z\to Y$ est
alors de type fini, et si $\mathcal H$ est le $\mathcal O_Z$-Module cohérent
tel que $u_*(\mathcal H)=\operatorname{gr}_{\mathcal J}^{\bullet}(\mathcal F)$,
il suffit d'appliquer à $u$ et à $\mathcal H$ le théorème de platitude
générique (6.9.1) pour prouver la proposition.
\end{proof}

\begin{proposition}[6.10.3]
\label{IV.2.s6.10.3-fr}
Les notations étant celles de (6.10.1), supposons que $\mathcal F$ soit
normalement plat le long de $Y$. Alors:
\begin{enumerate}
\item[(i)] $Y\cap\operatorname{Supp}(\mathcal F)$ est une partie à la fois
ouverte et fermée de $Y$ (autrement dit, une réunion de composantes connexes
de $Y$).
\item[(ii)] Si $\mathcal J$ est localement nilpotent,
$\operatorname{Supp}(\mathcal F)$ est une partie ouverte et fermée de $X$, et
pour tout $x\in\operatorname{Supp}(\mathcal F)$, on a
\end{enumerate}
\begin{align}
\dim(\mathcal F_x)&=\dim(\mathcal O_{Y,x})
  \tag{6.10.3.1}\label{IV.2.s6.10.3.1-fr}\\
\operatorname{prof}(\mathcal F_x)&=\operatorname{prof}(\mathcal O_{Y,x})
  \tag{6.10.3.2}\label{IV.2.s6.10.3.2-fr}\\
\operatorname{coprof}(\mathcal F_x)&=\operatorname{coprof}(\mathcal O_{Y,x}).
  \tag{6.10.3.3}\label{IV.2.s6.10.3.3-fr}
\end{align}
\end{proposition}

\begin{proof}
\begin{enumerate}
\item[(i)] Avec les notations de (6.10.1), on a
$Y\cap\operatorname{Supp}(\mathcal F)=\operatorname{Supp}(\mathcal G_0)$
(I, 9.1.13), et comme par hypothèse $\mathcal G_0$ est un
$\mathcal O_Y$-Module localement libre de type fini, son support est à la fois
ouvert et fermé dans $Y$.

